\documentclass[11pt,a4paper]{article}

\usepackage[margin=1in]{geometry}
\usepackage{microtype}
\usepackage{hyperref}
\usepackage{booktabs}
\usepackage{enumitem}
\usepackage{xcolor}
\usepackage{titlesec}
\usepackage{graphicx}
\usepackage{listings}
\usepackage[T1]{fontenc}
\usepackage{lmodern}

\definecolor{accent}{HTML}{2563EB}
\definecolor{codebg}{HTML}{F5F5F5}
\definecolor{codetext}{HTML}{333333}

\hypersetup{
  colorlinks=true,
  linkcolor=accent,
  urlcolor=accent,
}

\titleformat{\section}{\large\bfseries\color{accent}}{}{0em}{}
\titleformat{\subsection}{\normalsize\bfseries}{}{0em}{}

\lstset{
  backgroundcolor=\color{codebg},
  basicstyle=\small\ttfamily\color{codetext},
  breaklines=true,
  frame=single,
  rulecolor=\color{codebg},
  xleftmargin=1em,
  framexleftmargin=0.5em,
}

\title{\textbf{Greatwheel}\\[0.3em]
  \large An rLM-Coordinated Agentic Runtime:\\
  System Overview and Experimental Findings}
\author{M.\ Berends}
\date{March 2026}

\begin{document}
\maketitle

\begin{abstract}
Greatwheel is an agentic runtime built on the \emph{Recursive Language Model}
(rLM) paradigm: language models write Python code in sandboxed REPLs, using
variables---not the context window---as their working memory. A Rust core
manages execution, routing, and observability while Python agents run inside
\texttt{ouros} sandboxes. This report describes the system architecture,
summarises experiments conducted on the BrowseComp-Plus factoid retrieval
benchmark, and outlines planned next steps.
\end{abstract}

% -----------------------------------------------------------------
\section{Introduction}

Most agentic systems pass retrieved context directly into an LLM's context
window. Greatwheel takes a different approach: the model never sees full
context. Instead, every external operation---LLM calls, memory queries,
inter-agent messages---is a \emph{pause point} where the Rust runtime
intercepts, validates permissions, records traces, and resumes execution. The
agent's working state lives as Python variables that persist across turns,
giving the system perfect auditability and fine-grained resource control.

A ``bare agent'' with no predefined code is still functional: the rLM can
compose SDK globals (\texttt{llm}, \texttt{memory}, \texttt{agent},
\texttt{channel}) into ad-hoc tools on the fly, and those tools survive to the
next turn because they are ordinary Python functions in a persistent session.

% -----------------------------------------------------------------
\section{Architecture}

The system is organised as a nine-crate Rust workspace with the following
responsibilities:

\begin{description}[style=nextline, labelwidth=2.5cm, leftmargin=3cm]
  \item[\texttt{gw-core}] Shared types: IDs, Task, AgentDef, CallContext, permissions.
  \item[\texttt{gw-runtime}] \texttt{ouros} integration, REPL session management, host-function bridge.
  \item[\texttt{gw-llm}] LLM provider trait and Ollama client (via \texttt{rl-play} proxy).
  \item[\texttt{gw-memory}] Hybrid search: LanceDB vectors + tantivy BM25 + Postgres, with RRF fusion.
  \item[\texttt{gw-bus}] Inter-agent message bus (call/notify).
  \item[\texttt{gw-channels}] Channel adapters (HTTP, WebSocket, webhooks).
  \item[\texttt{gw-scheduler}] Task queue, per-org concurrency, rate limiting.
  \item[\texttt{gw-trace}] OpenTelemetry GenAI instrumentation with Postgres export.
  \item[\texttt{gw-server}] Binary crate: HTTP/WS API that wires everything together.
\end{description}

\noindent
The production stack runs four Docker Compose services: the Greatwheel binary,
\texttt{rl-play} (an LLM proxy that captures conversations for fine-tuning),
Postgres, and Ollama.

\subsection{The rLM Execution Loop}

\begin{enumerate}
  \item A user message creates a \textbf{Task} dispatched to an agent's
        \texttt{ouros} session.
  \item Context is injected as a Python variable
        (\texttt{\_task = \{"payload": ...\}}).
  \item The rLM receives a system prompt, a truncated session summary, and the
        user message. It outputs Python code.
  \item \texttt{ouros} executes the code. Pure computation runs immediately;
        host-function calls (\texttt{llm}, \texttt{memory}, etc.) \emph{pause}
        the session, and the Rust runtime fulfils each call before resuming.
  \item The loop repeats until the rLM produces a \texttt{FINAL()} response or
        an iteration limit is reached.
  \item All Python variables persist to the next conversational turn.
\end{enumerate}

\subsection{Observability and Fine-Tuning}

Every operation emits OpenTelemetry spans following GenAI semantic conventions.
Spans are exported to the console, OTLP, and Postgres (the latter enables
fine-tuning dataset extraction). All LLM calls also flow through
\texttt{rl-play}, which groups conversations by session and supports SFT, DPO,
and PPO exports---creating a closed loop from production traces to improved
models.

% -----------------------------------------------------------------
\section{BrowseComp-Plus Experiments}

BrowseComp-Plus is a factoid retrieval benchmark: 30 queries over a 100K web
document corpus. It serves as the primary testbed for our rLM + retrieval
research. We have conducted over 40 experiment runs, methodically varying
retrieval strategy, prompting, and model configuration.

\subsection{Current Best Result}

Our best configuration achieves \textbf{36.7\% (11/30)} with a reproducibility
range of 9--11 across runs (mean $\approx$10). Per-query timing averages 115\,s,
of which 90\% is LLM inference in the rLM loop and 10\% is pre-search
(decomposition + BM25 + PRF).

\subsection{What Has Worked}

\paragraph{Retrieval.}
\begin{itemize}[nosep]
  \item \textbf{Boosted BM25}: Phrase match (4$\times$), slop phrase (2$\times$),
        AND'd terms (1.5$\times$), OR'd terms (1$\times$). This was the single
        largest contributor, lifting accuracy by 2--3 queries.
  \item \textbf{Pseudo-relevance feedback (PRF)}: Extracting distinctive terms
        from round-1 snippets for expanded queries added $\sim$1--2 queries.
  \item \textbf{Multi-query decomposition}: Breaking multi-hop questions into
        independent sub-fact searches (1$\to$5 queries) meaningfully broadened
        retrieval coverage.
  \item \textbf{$k=10$ retrieval breadth}: More documents per search without
        overwhelming the model.
\end{itemize}

\paragraph{Prompting and loop control.}
\begin{itemize}[nosep]
  \item \textbf{Document-grounded prompts}: Forcing \texttt{get\_document()} +
        \texttt{llm\_query()} and emphasising corpus-only answers.
  \item \textbf{VERIFY directive}: ``Search for the candidate answer by name to
        confirm'' catches extraction errors.
  \item \textbf{12 iterations}: Enough to explore, verify, and submit. 16
        iterations caused second-guessing.
  \item \textbf{\texttt{think:false} for utility calls}: Saves tokens on
        pre-search extraction and fallback.
  \item \textbf{Fallback extraction}: Recovers answers when the model does not
        call \texttt{FINAL()} in time.
\end{itemize}

\subsection{What Has Not Worked}

\begin{itemize}[nosep]
  \item \textbf{HyDE (hypothetical document embeddings)}: No benefit for
        BM25 queries.
  \item \textbf{Dense vector search} (all six variants tested): Added zero
        lift to BM25-only. Variants included asymmetric prefixes, BM25-weighted
        RRF fusion, increased IVF-PQ probes, and a fully re-embedded corpus.
        BrowseComp queries are entity-heavy (names, dates, places) where BM25
        keyword matching naturally excels; 768-dimensional dense vectors
        (\texttt{nomic-embed-text}) lack discrimination at 100K scale.
  \item \textbf{Best-of-$N$ majority voting}: Wrong answers agreed more often
        than correct ones---3$\times$ cost, same accuracy.
  \item \textbf{Wider pre-search} (8+5 queries): Overwhelmed the model with
        context.
\end{itemize}

\subsection{Key Insight: The Retrieval Ceiling}

\textbf{86\% of failures are retrieval misses}---the relevant documents are
never found. The 9B model reasons well once it has the right documents; the
bottleneck is search, not extraction. However, the missing documents often lack
both keyword \emph{and} semantic overlap with the query. They require reasoning
about \emph{what kind of document would contain the answer}, which neither BM25
nor dense retrieval performs.

% -----------------------------------------------------------------
\section{Challenges}

Beyond the retrieval ceiling, several broader challenges have shaped
development:

\begin{enumerate}[nosep]
  \item \textbf{High per-query variance.} The same configuration scores 9--11
        out of 30 across runs. Individual queries flip between correct and
        incorrect nondeterministically, making it hard to distinguish real gains
        from noise.
  \item \textbf{Token cost.} Each query consumes $\sim$150K tokens across a
        12-iteration loop. The rLM loop dominates wall-clock time, and reducing
        iterations hurts accuracy.
  \item \textbf{Incomplete runtime.} Several crates (\texttt{gw-bus},
        \texttt{gw-scheduler}, \texttt{gw-channels}) have trait definitions but
        lack concrete implementations. The Python-side agent SDK is designed but
        not yet wired to the host bridge.
  \item \textbf{Multi-tenancy.} Core types support org/user/session hierarchy
        and per-agent permissions, but auth middleware and session lifecycle
        (idle timeout, snapshot, eviction) are not yet implemented.
\end{enumerate}

% -----------------------------------------------------------------
\section{Next Steps}

Our roadmap separates retrieval-quality improvements from orchestration and
infrastructure work.

\subsection{Near-Term: Retrieval}

\begin{description}[style=nextline, labelwidth=4.5cm, leftmargin=5cm]
  \item[ColBERT reranking (R2)] Use Reason-ModernColBERT (149M params) as
    either a primary retriever or a BM25 reranker. ColBERT's late-interaction
    token-level MaxSim scoring has demonstrated strong results on
    BrowseComp-Plus in external work. We plan three modes: ColBERT-only, BM25 +
    ColBERT RRF fusion, and BM25$\to$ColBERT rerank. Requires $\sim$10--16\,GB
    VRAM.
  \item[Larger rLM model (A3)] Use a 32B model (\texttt{qwen2.5:32b} or
    \texttt{qwen3.5:32b}) for the main rLM loop, keeping the 9B model for
    utility calls. Better reasoning should produce better search queries and
    extraction.
  \item[Ensemble (A4)] Run 2--3 diverse configurations per query and take the
    union of correct answers. Historical analysis suggests the union across
    existing configs reaches $\sim$40\%.
\end{description}

\subsection{Near-Term: Optimisation Infrastructure}

\begin{description}[style=nextline, labelwidth=4.5cm, leftmargin=5cm]
  \item[GEPA integration] Use GEPA's \texttt{optimize\_anything} API for
    systematic prompt and parameter optimisation. The key unlock is
    \emph{actionable side information} (ASI): per-query failure-mode
    diagnostics that tell the optimiser \emph{why} a run scored as it did, not
    just the scalar accuracy.
  \item[ASI infrastructure] Even without GEPA, structured per-query diagnostics
    (failure mode, retrieval recall, searches issued, tokens consumed) enable
    smarter manual iteration.
\end{description}

\subsection{Medium-Term: Runtime Completion}

\begin{itemize}[nosep]
  \item \textbf{Episodes}: Compressed summaries of agent work that flow as
        context between agents, enabling composability and multi-agent
        collaboration.
  \item \textbf{Agent SDK}: Wire Python-side \texttt{llm}, \texttt{memory},
        \texttt{agent}, \texttt{channel} globals to the host bridge.
  \item \textbf{Session lifecycle}: Idle timeout, snapshot to Postgres, eviction.
  \item \textbf{Scheduler and bus}: Concrete implementations of task queuing,
        rate limiting, and inter-agent messaging.
\end{itemize}

% -----------------------------------------------------------------
\section{Conclusion}

The core architectural bet---context as Python variables, not context
window---has proven sound. The rLM + ouros execution model delivers
auditability, composability, and resource control that would be difficult to
achieve in a purely prompt-based system. On BrowseComp-Plus, boosted BM25 with
pseudo-relevance feedback and document-grounded prompting reaches 37\%
accuracy, with clear evidence that the retrieval ceiling (not model reasoning)
is the binding constraint. ColBERT reranking and larger models are the most
promising levers for the next phase of work.

\end{document}
